% Esquema de la propuesta metodológica (etapas de CRISP-DM y fases del PFI)
\begin{tikzpicture}[
    node distance = 0.4cm,
    etapa/.style = {draw, align=center, font=\small, fill=gray!10,
                    text width=2.05cm, minimum height=1.3cm, inner sep=3pt,
                    execute at begin node={\setstretch{1}\hyphenpenalty=10000}},
    fase/.style = {font=\footnotesize, align=center,
                   execute at begin node={\setstretch{1}}},
    llave/.style = {decorate, decoration={brace, amplitude=5pt, mirror}, thick},
    flecha/.style = {-{Stealth[length=2.2mm]}, thick}
  ]
  \node[etapa] (e1) {Comprensión del problema};
  \node[etapa, right=of e1] (e2) {Comprensión de los datos};
  \node[etapa, right=of e2] (e3) {Preparación de los datos};
  \node[etapa, right=of e3] (e4) {Modelado};
  \node[etapa, right=of e4] (e5) {Evaluación};
  \node[etapa, right=of e5, fill=gray!25] (e6) {Despliegue};

  \foreach \a/\b in {e1/e2, e2/e3, e3/e4, e4/e5, e5/e6}
    \draw[flecha] (\a) -- (\b);

  % Iteraciones de CRISP-DM
  \draw[flecha] (e5.north) -- ++(0,0.45cm) -| node[fase, pos=0.25, above] {nuevas iteraciones} (e3.north);

  % Fases del PFI
  \draw[llave] ($(e1.south west)+(0,-0.15cm)$) -- node[fase, below=6pt] {Primera fase:\\comprensión y selección de datos} ($(e2.south east)+(0,-0.15cm)$);
  \draw[llave] ($(e3.south west)+(0,-0.15cm)$) -- node[fase, below=6pt] {Segunda fase:\\modelado} ($(e5.south east)+(0,-0.15cm)$);
  \draw[llave] ($(e6.south west)+(0,-0.15cm)$) -- node[fase, below=6pt] {Tercera fase:\\aplicación web} ($(e6.south east)+(0,-0.15cm)$);
\end{tikzpicture}
